
%----------------------------------------
\begin{frame}[fragile]
  \frametitle{Commands and tests} \small

\begin{definition}[Command]
A command is a dictionary mapping command name and arguments:

\begin{lstlisting}[language=python,frame=single]
Command = Dict[str, Union[int, float, str]]
\end{lstlisting}
\end{definition}

Example of command $goto(x=100,y=200)$:

\[
["name" : "goto", "x": 100, "y": 200]
\]

\pause
\begin{definition}[Test]
A test is a list of commands

\begin{lstlisting}[language=python,frame=single]
Test = list[Command]
\end{lstlisting}
\end{definition}

\end{frame}

%----------------------------------------
\begin{frame}<0>{Hidden slide}
  \frametitle{Matching Temporal Logic (MaTL) - the match construct} \small

Consider the trace:

\[ 
[ 
open(file=\texttt{"log"}, mode=\texttt{"r"}),
read(file=\texttt{"log"}),
close(file=\texttt{"log"})
]
\]

\pause

\begin{block}{We want a formula of the form}
\begin{equation*}
\begin{aligned}
\varphi \;::=&\; ...  \;|\; \langle id(m) \rangle \varphi \;|\; ...
\end{aligned}
\end{equation*}
\end{block}

\pause

\vspace{0.2cm}

For example, at the first trace position the following formula is true,
\underline{binding $f$ to $\texttt{"log"}$},
if $\varphi$ thereafter is true in scope of this binding:

\[
\langle open(file=f?, mode="r") \rangle \varphi
\]

\pause

\vspace{0.2cm}

At the last position the following formula is true, \underline{assuming $f$ is in scope} and denotes $\texttt{"log"}$), if $\varphi$ is thereafter true:

\[
\langle close(file=f) \rangle \varphi
\]

\end{frame}

%----------------------------------------

\begin{frame}<0>
  \frametitle{A test is a list of commands} \small

\begin{block}{Property}\scriptsize
If at some position in the test there is a MOVE command with a distance `$d$’, then later there must be a ROTATE command
with an angle `$a$’ such that the predicate $P(d,a)$ holds.
\end{block}

\begin{figure}
\includegraphics[width=1.0\textwidth]{figures/slide1.png}
%\caption{Your caption}
\end{figure}

\begin{block}{}\scriptsize
A test is a sequence of commands. A test
can satisfy a temporal property or not.
\end{block}

\end{frame}

%----------------------------------------
\begin{frame}<0>
  \frametitle{MaTL syntax} \small

\begin{equation*}
\begin{aligned}
\varphi \;::=&\; \true 
  \;|\; \neg \varphi 
  \;|\; \varphi_1 \lor \varphi_2 \pause
  \;|\; \langle id(m) \rangle \varphi \pause
  \;|\; e_1 \; \diamond \; e_2 
  \;|\; e \vdash r \\[0.5ex] \pause
    & \Next \varphi 
  \;|\; \varphi_1 \until \varphi_2 
  \;|\; \prev \varphi 
  \;|\; \varphi_1 \since \varphi_2 \pause
  \;|\; \Sigma^{+}_{[a,b]} \varphi 
  \;|\; \Sigma^{-}_{[a,b]} \varphi \\[0.5ex] \pause \\
m \;::=&\; \epsilon \;|\; m_1 \; m_2 \;|\; id = p \\[0.5ex]
p \;::=&\; x? \;|\; x \;|\; c \;|\; id_1.id_2 \\[0.5ex] \pause \\
e \;::=&\; x \;|\; c \;|\; id_1.id_2 \;|\; e_1 \; \otimes \; e_2 \\[0.5ex]
\diamond \;::=&\; < \;|\; \le \;|\; = \;|\; \neq \;|\; \ge \;|\; > \\[0.5ex]
\otimes \;::=&\; + \;|\; - \;|\; \times \;|\; \div 
\end{aligned}
\end{equation*}

\end{frame}


%----------------------------------------
\begin{frame}<0>
  \frametitle{Example} \small

\begin{block}{Requirement}
Whenever a file f is opened, it is eventually closed.
\end{block}

What we can write using the core syntax:

\[
\neg (\true \until (\langle open(file=f?)\rangle \neg(\true \until close(file=f))))
\]

\pause

We might instead want to to write:

%\[
%\Box ( [ open(file=f?) ] \Diamond close(file=f))
%\]

%\pause

%Or even:

\[
\Box (open(file=f?) \Rightarrow \Diamond close(file=f))
\]

\vspace{0.2cm}

% The latter two forms are derived syntax.

The last form is {\bf derived} syntax.

\end{frame}


%----------------------------------------
\begin{frame}<0>
  \frametitle{Derived syntax} \small

\[
\begin{array}{rcl@{\hspace{2em}}ll}

\false & = & \neg \true                     & \text{false} & (1)\\

\varphi_1 \wedge \varphi_2 &=& \neg(\neg \varphi_1 \vee \neg \varphi_2) & \text{and} & (2)  \\

\varphi_1 \rightarrow \varphi_2 &=& \neg \varphi_1 \vee \varphi_2       & \text{implication} & (3) \\

%e_1 \diamond_1 e_2 \diamond_2 e_3 &=& e_1 \diamond_1 e_2 \wedge e_2 \diamond_2 e_3 & \text{ternary relation} & (4)\\ 

\\

%\pause

[id(m)] \varphi &=& \neg \langle id(m) \rangle \neg \varphi & \text{universal modality} & (5)\\ 

%\pause

%id(m) \Rightarrow  \varphi &=& [id(m)] \varphi & \text{universal modality} & (6)\\

%\pause

%id(m) \bullet \varphi &=& \langle id(m) \rangle \varphi & \text{existential modality} & (7)\\

%\pause

id(m) &=& \langle id(m) \rangle \true      & \text{event occurrence} & (8)\\

\\

%\pause

\Diamond \varphi &=& \true \until \varphi & \text{eventually in the future} & (9)\\

\Box \varphi &=& \neg \Diamond \neg \varphi & \text{always in the future} & (10)\\

\DDiamond \varphi &=& \true \since \varphi & \text{sometime in the past} & (11)\\

\BBox \varphi &=& \neg \DDiamond \neg \varphi & \text{always in the past} & (12)\\ \\

\varphi_1 \wuntil \varphi_2 &=& (\varphi_1 \until \varphi_2) \vee \Box \varphi_1 & \text{weak until} & (13)\\

\varphi_1 \wsince \varphi_2 &=& (\varphi_1 \since \varphi_2) \vee \BBox \varphi_1 & \text{weak since} & (14)\\

\wnext \varphi &=& \next \varphi \vee \neg \next \true & \text{weak next} & (15)\\

\wprev \varphi &=& \prev \varphi \vee \neg \prev \true & \text{weak previous} & (16)\\

\end{array}
\]

\end{frame}


%----------------------------------------
\begin{frame}<0>
  \frametitle{Derived syntax} \small

\[
\begin{array}{rcl@{\hspace{2em}}ll}

\Sigma^{+}_{a} \varphi  &=& \Sigma^{+}_{[a,a]} \varphi & \text{future exact count} & (17)\\

\Sigma^{-}_{a} \varphi &=& \Sigma^{-}_{[a,a]} \varphi & \text{past exact count} & (18)\\

%\pause

\Next_a \varphi &=& \Next \Next ... \Next  \varphi  \quad (\text{a times}) & \text{a steps in future} & (19)\\

\prev_a \varphi &=& \prev \prev ... \prev \varphi \quad (\text{a times}) & \text{a steps in past} & (20)\\

\end{array}
\]

\end{frame}


%----------------------------------------
\begin{frame}<0>
  \frametitle{A small file example} \small

An example trace:

\[ 
[
open(file=\texttt{"log"}, mode=\texttt{"r"}),
read(file=\texttt{"log"}),
close(file=\texttt{"log"})
]
\]

\pause

\begin{block}{Requirement}
(1) Whenever a file $f$ is opened, it is eventually closed.\\
(2) Whenever a file $f$ is read from, in the past it has been opened in read mode, and not closed since.\\
(3) There should be one open event in read mode.
\end{block}

\pause

\begin{equation*}
\begin{aligned}
(1) \;&\; \Box (open(file=f?) \Rightarrow \Diamond close(file=f)) \\[0.5ex]
(2) \;&\; \Box (read(file=f?) \Rightarrow \neg close(file=f) \since open(file=f, mode=\texttt{"r"}) \\[0.5ex]
(3) \;&\; \Sigma^{+}_1 open(mode=\texttt{"r"})
\end{aligned}
\end{equation*}

\end{frame}

